% !TeX encoding = UTF-8
\documentclass[variant=apuntes,cover=institutional,mode=screen]{ua-engineering-v6-article}
% !TeX encoding = UTF-8
\UASetUniversity{Universidad Austral}
\UASetFaculty{Facultad de Ingeniería}
\UASetDepartment{Departamento de Computación}
\UASetCampus{Pilar}
\UASetCareer{Ingeniería Informática}
\UASetCommission{Comisión A}
\UASetProfessor{Equipo docente}
\UASetSemester{Segundo cuatrimestre 2026}
\UASetCourse{Sistemas Distribuidos}
\UASetDate{2026}


\UASetClassNumber{03}
\UASetClassTitle{Consistencia y replicación}
\UASetDocKind{Apuntes}

\title{Clase 03 --- Consistencia, replicación y modelos de estado}
\author{\UAProfessor}
\date{\UADate}

\begin{document}
\maketitle

\section{Introducción}

Una vez resuelta la comunicación (RPC), el problema central es el \textbf{estado}. Para mejorar disponibilidad y latencia, los sistemas replican datos. Esto introduce una tensión fundamental:

\begin{quote}
No es posible mantener múltiples copias perfectamente sincronizadas sin coordinación.
\end{quote}

\section{Modelo}

Consideremos un objeto replicado en $N$ nodos. Las operaciones son:
\begin{itemize}
\item \texttt{write(k, v)}
\item \texttt{read(k)}
\end{itemize}

La red es asíncrona: mensajes pueden retrasarse, perderse o reordenarse.

\section{Divergencia}

Dos réplicas $R_i$ y $R_j$ divergen si almacenan valores distintos para la misma clave en el mismo tiempo lógico. Esto ocurre por:
\begin{itemize}
\item latencias distintas
\item concurrencia
\item particiones
\end{itemize}

\section{Consistencia}

\subsection{Linearizabilidad}

Una historia es linearizable si existe un orden total que:
\begin{itemize}
\item respeta el tiempo real
\item es consistente con la especificación secuencial
\end{itemize}

\paragraph{Propiedad}
Cada operación tiene un punto de linearización entre invocación y respuesta.

\subsection{Consistencia causal}

Definimos la relación \emph{happens-before} ($\rightarrow$):
\begin{itemize}
\item orden de programa
\item envío/recepción de mensajes
\item transitividad entre relaciones conocidas
\end{itemize}

La consistencia causal exige que si $a \rightarrow b$, entonces todo proceso observa $a$ antes que $b$.

\section{Tiempo físico, orden lógico y causalidad}

En sistemas distribuidos no existe un reloj global perfecto que permita ordenar todos los eventos con certeza. Cada nodo observa su propio reloj local, la red introduce demoras variables y los mensajes pueden llegar tarde, perderse o reordenarse. Por eso, muchas decisiones de diseño deben basarse en causalidad lógica y no solamente en timestamps físicos.

\begin{uakeybox}
La pregunta importante no siempre es ``qué hora marcaba el reloj'', sino ``qué evento pudo influir sobre qué otro evento''.
\end{uakeybox}

Comparar timestamps físicos entre nodos distintos puede ser engañoso:
\begin{itemize}
\item los relojes locales pueden derivar;
\item una corrección de sincronización puede mover un reloj hacia adelante o hacia atrás;
\item un mensaje enviado antes puede llegar después por latencia variable;
\item dos eventos pueden tener timestamps comparables sin que exista relación causal entre ellos.
\end{itemize}

\subsection{Relojes de Lamport}

Los relojes de Lamport asignan un contador lógico a cada evento. Cada proceso mantiene un contador $L_i$:
\begin{itemize}
\item antes de un evento local, incrementa $L_i$;
\item al enviar un mensaje, adjunta el valor actual;
\item al recibir un mensaje con timestamp $t$, actualiza $L_i \leftarrow \max(L_i,t)+1$.
\end{itemize}

La propiedad central es:
\[
a \rightarrow b \implies L(a) < L(b)
\]

La implicación inversa no vale: si $L(a) < L(b)$, no necesariamente $a \rightarrow b$. Lamport preserva causalidad conocida, pero no detecta concurrencia con precisión.

\subsection{Consistencia eventual}

Garantiza convergencia:
\begin{quote}
Si no hay nuevas escrituras, todas las réplicas convergen al mismo valor.
\end{quote}

No garantiza orden ni lecturas recientes.

\section{Arquitecturas de replicación}

\subsection{Leader-based}

Un líder serializa escrituras. Replicación:
\begin{itemize}
\item síncrona: mayor consistencia, más latencia
\item asíncrona: menor latencia, riesgo de pérdida
\end{itemize}

\subsection{Multi-leader}

Cada región acepta escrituras. Problema:
\begin{itemize}
\item conflictos concurrentes
\end{itemize}

\subsection{Leaderless (Dynamo)}

Parámetros:
\begin{itemize}
\item $N$: número de réplicas
\item $W$: write quorum
\item $R$: read quorum
\end{itemize}

\paragraph{Intersección}
Si $R+W>N$, los quorums se intersectan.

\section{Conflictos}

Dos escrituras concurrentes generan versiones divergentes.

\subsection{Versionado}

Vector clocks:
\begin{itemize}
\item detectan concurrencia
\item permiten merge explícito
\end{itemize}

Cada proceso mantiene un vector con una componente por proceso. Para comparar dos vectores $V(a)$ y $V(b)$:
\begin{itemize}
\item $V(a) < V(b)$ si todas las componentes de $V(a)$ son menores o iguales que las de $V(b)$ y al menos una es estrictamente menor;
\item si $V(a) < V(b)$, entonces $a$ precede causalmente a $b$;
\item si los vectores no son comparables, los eventos son concurrentes.
\end{itemize}

Por ejemplo, $(2,1,0)$ y $(1,2,0)$ no son comparables; representan escrituras concurrentes. En un store replicado, esto indica que no conviene sobrescribir ciegamente una versión con otra, sino conservar ambas o aplicar reconciliación semántica.

\subsection{Resolución}

\begin{itemize}
\item LWW: simple, puede perder datos
\item merge semántico
\end{itemize}

\section{Anti-entropy}

Mecanismos:
\begin{itemize}
\item gossip
\item read-repair
\item hinted handoff
\end{itemize}

\section{Trade-offs}

\begin{itemize}
\item consistencia fuerte → coordinación → latencia
\item consistencia débil → disponibilidad → anomalías
\end{itemize}

\section{Marco $D=(P,F,G,S,T)$}

\begin{itemize}
\item $P$: mantener estado replicado
\item $F$: delay, pérdida, partición
\item $G$: modelo de consistencia
\item $S$: arquitectura (leader, quorum)
\item $T$: latencia, disponibilidad, complejidad
\end{itemize}

\section{Conclusión}

\begin{quote}
La consistencia no es una propiedad binaria: es un contrato explícito entre el sistema y sus clientes.
\end{quote}


\section{Síntesis razonada de la clase}
Esta unidad debe leerse como parte de una cadena de decisiones de diseño y no como un catálogo de mecanismos. Para estudiar el tema con profundidad conviene reconstruir cuatro preguntas: \emph{qué problema aparece}, \emph{qué garantía se desea}, \emph{qué mecanismo la aproxima} y \emph{qué costo introduce}. En cada ejemplo debe identificarse además el modelo de falla implícito.

\begin{uakeybox}
Una respuesta técnicamente madura no termina al nombrar un algoritmo o patrón: declara sus supuestos, la propiedad que preserva y el escenario en el que deja de ser suficiente.
\end{uakeybox}

\section{Bibliografía guiada y ruta de profundización}
Para profundizar esta clase se recomienda: Kleppmann, caps. 5, 8 y 9; Lamport sobre tiempo y orden.

La lectura debe hacerse con preguntas concretas: (1) ¿qué modelo de sistema asume el autor?, (2) ¿qué propiedad demuestra o persigue?, (3) ¿qué falla o anomalía motiva el mecanismo?, y (4) ¿qué trade-off queda abierto?

\section{Conexión con el Trabajo Final Integrador}
El grupo debe señalar qué parte de su sistema utiliza los conceptos de esta unidad, justificar por qué son necesarios y documentar una alternativa descartada. Cuando el contenido todavía no corresponda a la etapa vigente, debe registrarse como hipótesis o decisión pendiente.

\section{Autoevaluación conceptual}
\begin{enumerate}
\item ¿Cuál es el problema distribuido concreto que motiva esta clase?
\item ¿Qué propiedad de seguridad y qué propiedad de progreso están en juego?
\item ¿Qué supuesto de red, tiempo o falla resulta decisivo?
\item ¿Qué trade-off aceptaría en un sistema real y cuál no aceptaría en un dominio crítico?
\item ¿Cómo observaría experimentalmente que la solución se comporta como espera?
\end{enumerate}

\section{Profundización autocontenida v9}
\subsection{Problema, límite e invariante}
El núcleo de esta clase es replicación síncrona y asíncrona, stale reads, quórums, causalidad y modelos de consistencia. Para estudiarlo de manera autónoma conviene separar tres planos. Primero, el \emph{problema observable}: qué comportamiento incorrecto puede ver un cliente o un operador. Segundo, el \emph{límite del modelo}: qué información, sincronía o coordinación no está disponible. Tercero, el \emph{invariante}: qué propiedad no debe violarse aunque el sistema pierda progreso temporalmente.

El modelo de fallas relevante incluye réplica atrasada, escritura concurrente, partición y failover. Una solución debe describirse como protocolo y no solo como nombre de tecnología: actores, estados, mensajes, condición de éxito, condición de aborto y estado visible después de una interrupción. La evidencia esperada es historia de lecturas/escrituras, staleness y comparación de quórums.

\subsection{Método de análisis}
Ante un caso nuevo, siga esta secuencia:
\begin{enumerate}
\item Identifique actores, dato o recurso compartido y frontera de confianza.
\item Escriba una ejecución nominal y una ejecución adversarial.
\item Distinga seguridad (lo malo no ocurre) de progreso (algo bueno finalmente ocurre).
\item Declare qué supuesto permite avanzar: mayoría, deadline, sincronía parcial, operación idempotente, almacenamiento durable u otro.
\item Compare una alternativa más simple y una más fuerte; registre qué métrica o experimento haría cambiar la decisión.
\end{enumerate}

\subsection{Caso transferible}
Considere una plataforma de reservas que recibe operaciones duplicadas, pierde conectividad entre regiones y debe sostener un camino crítico sin ocultar el estado real al usuario. Aplique el mecanismo de esta clase a una única decisión concreta. Una respuesta madura no intenta resolver todo: delimita la operación, formula la garantía y explica la degradación esperada cuando el supuesto central deja de cumplirse.

\subsection{Errores de razonamiento frecuentes}
\begin{itemize}
\item confundir una propiedad con el producto que suele implementarla;
\item describir el camino feliz sin recorrer una falla intermedia;
\item afirmar ``exactly-once'', ``alta disponibilidad'' o ``consistencia'' sin definir alcance observable;
\item medir solo throughput aunque la hipótesis trate sobre semántica o recuperación;
\item presentar la complejidad añadida como beneficio sin compararla con la alternativa más simple.
\end{itemize}

\section{Lectura guiada}
Kleppmann, capítulo 5; lecturas sobre consistencia causal y quórums.

Durante la lectura, registre: modelo de sistema, propiedad perseguida, contraejemplo que motiva el mecanismo, costo principal y una condición bajo la cual no lo elegiría. Estas cinco anotaciones convierten la bibliografía en una herramienta de diseño y no en una lista para memorizar.

\section{Aplicación al Trabajo Final Integrador}
El equipo debe producir un ADR breve con la decisión asociada a esta clase. Debe contener: contexto, dos alternativas, supuesto decisivo, garantía elegida, trade-off, evidencia pendiente y fecha de revisión. La evidencia mínima esperada para esta unidad es: historia de lecturas/escrituras, staleness y comparación de quórums.

\section{Autoevaluación avanzada}
\begin{enumerate}
\item ¿Puedo explicar la clase sin nombrar primero una tecnología?
\item ¿Puedo construir una ejecución que viole la garantía si el mecanismo se configura mal?
\item ¿Puedo distinguir seguridad, progreso y experiencia del usuario?
\item ¿Puedo proponer una medición que valide la propiedad, no solo el rendimiento?
\item ¿Puedo defender cuándo una solución más simple sería suficiente?
\end{enumerate}

\end{document}
